\section{拒稿诊断与本计划的任务}

IP\&M 拒稿意见的第二句是：所报告的研究结果过于不成熟（the research results reported are too premature for publication），并要求补充更多工作以支撑论文结论（more work is needed to substantiate the conclusions）。与第一句（新颖性）不同，这一意见指向证据链本身，而且它在论文的局限一节里可以逐条对上号。把局限一节的五条边界与拒稿意见对照阅读，可以精确拆出``premature''的三个具体构成。

其一，端到端证据的规模失衡。选择器层面的评估是 1{,}600 条记录、六个选择器、分离的配对区间——这一层是成熟的；但端到端问答只有单一阅读器（GLM-4-Plus）、单一数据集（QASPER）、单一操作点（$4\times$）、$n{=}25$，其配对 F1 差的 95\% 区间为 $[-14.5, 8.2]$，半宽超过 $\pm 11$ 个 F1 点。任何关于端到端效应的实质性主张（例如``\sieve{} 与截断统计持平''）在这个区间宽度下都只能是方向性陈述，而编辑与审稿人读到的是``前层实验充分、终层实验示踪''的观感。其二，机制消融未连成矩阵。组件消融（去问题条件化、平位置先验、去冗余惩罚、纯相关性）目前只在 QASPER 单数据集上运行，且与预算扫描、位置研究、效率测量各自成节，没有形成``预算 $\times$ 配置 $\times$ 指标 $\times$ 数据集''的统一矩阵，读者无法直接看到三个选择目标在整个预算范围内的交换率。其三，结论的单点依赖。论文的核心叙事——不同选择目标产生结构性权衡——目前依赖单一阅读器的间接证据（召回--F1 相关性 $n{=}25$/方法），缺少跨阅读器的稳健性背书。

本计划给出三条互补的补强路径，分别对准上述三个构成：路径 A 扩大端到端评估（最高优先级，直接回应``premature''）；路径 B 引入阅读器多样性并重测阅读器侧位置剖面（回应单点依赖与借用的位置剖面）；路径 C 把消融扩展成跨数据集统一矩阵（纯本地计算、零 API 依赖、成本最低）。三条路径与第一阶段的 novelty 重构严格对齐：第一阶段的结论是论文的新颖性由``formulation + frontier 方法论 + 结构性权衡发现''三根支柱承担，因此实验补强的目标不是让 \sieve{} 在任何指标上击败所有基线——那既不可能也不必要——而是让三根支柱的每一句声明都有相称规模的证据。负结果（例如 BM25 召回更高）在新的矩阵里应被更清晰地呈现为权衡点而非失败，这正是实验设计与写作纪律的结合处，第 \ref{sec:discipline} 节单独规定。

执行顺序遵循两个原则。第一，先便宜后昂贵：路径 C 全部为本地 CPU 计算（仓库现有 M 系列脚本的小改），路径 A 第一步是恢复执行已实现并发布的 $n{=}100$ 协议（脚本已带断点续跑），只有 A2/A3 与路径 B 需要新的 API 预算或 GPU。第二，先支撑后扩展：所有实验先在既有数据集与操作点上把声明的置信区间收紧，再考虑新数据集、新阅读器等增量维度；每一步的完成定义（第 \ref{sec:dod} 节）都映射回拒稿意见的具体条目，避免``多跑实验''式的盲目扩容。
